%------------------------------------------------------------------------------%

\usepackage{integracao_e_series}
\usepackage{ulem}

\title{Integração por Substituição Trigonométrica}

%------------------------------------------------------------------------------%
\begin{document}

\addtitle

%------------------------------------------------------------------------------%
\section{Integração por Substituição Trigonométrica}

%------------------------------------------------------------------------------%
\begin{frame}{Substituição Simples}
  Calculando
  \[
    F = \int x \sqrt{a^2-x^2}\,dx
  \]
  \pause
  substituição
  \(
    u = a^2-x^2
    \qquad\qquad
    du = -2xdx
    \qquad
    xdx = -\frac{du}{2}
  \)
  \pause
  \[
    F = -\frac{1}{2}\int \sqrt{u}\,du
  \]
\end{frame}

%------------------------------------------------------------------------------%
\begin{frame}{Substituição Simples -- Problema}
  Calculando
  \[
    F = \int \sqrt{a^2-x^2}\,dx
  \]
  \pause
  substituição
  \(
    u = a^2-x^2
    \qquad\qquad
    du = -2xdx
    \qquad
    dx = -\frac{du}{2x}
  \)
  \pause
  \[
    F
    = -\frac{1}{2}\int \sqrt{u} \frac{du}{x}
    \pause
    =\; ? \vphantom{\int}
  \]
\end{frame}

%------------------------------------------------------------------------------%
\begin{frame}{Segunda Tentativa}
  Calculando
  \[
    F = \int \sqrt{a^2-x^2}\,dx
  \]
  \pause
  substituição
  \(
    x = a\sin(\theta)
    \qquad\qquad
    dx = a\cos(\theta)d\theta
    \vphantom{\frac{du}{2}}
  \)
  \pause
  \[
    F = \int \sqrt{a^2 - \big(a\sin(\theta)\big)^2\,} \,a\cos(\theta)d\theta
  \]
\end{frame}

%------------------------------------------------------------------------------%
\begin{frame}{Segunda Tentativa}
  \begin{align*}
    \onslide<+->{F & = \int \sqrt{a^2 - \big(a\sin(\theta)\big)^2\,} \,a\cos(\theta)d\theta     \\[2mm]}
    \onslide<+->{  & = \int \sqrt{a^2 \big( 1 - \sin^2(\theta)\big)\,} \, a\cos(\theta) d\theta \\[2mm]}
    \onslide<+->{  & = \int a^2 \sqrt{ \cos^2(\theta) } \,\cos(\theta) d\theta}
    \onslide<+->{  & & a > 0        \\[2mm]}
    \onslide<+->{  & = \int a^2 \cos^2(\theta) d\theta
                   & & \theta \in \left[-\frac{\pi}{2},\, \frac{\pi}{2}\right] \\[2mm]}
    \onslide<+->  {& = a^2 \int \cos^2(\theta) d\theta}
  \end{align*}
\end{frame}

%------------------------------------------------------------------------------%
\begin{frame}{Segunda Tentativa}
  \begin{align*}
    \onslide<+->{F & = a^2 \int \cos^2(\theta) d\theta         \\[2mm]}
    \onslide<+->{F & = \frac{a^2}{2} \int \big(1+\cos(2\theta)\big) d\theta
      & & \cos^2(\alpha)= \frac{1}{2}\big(1+\cos(2\alpha)\big) \\[2mm]}
    \onslide<+->{  & = \frac{a^2}{2}\left(\theta + \frac{\sin(2\theta)}{2}\right) + C \\[2mm]}
    \onslide<+->{  & = \frac{a^2}{2}\big(\theta + \sin(\theta)\cos(\theta)\big) + C \\[2mm]}
    \onslide<+->{  & = \frac{a^2}{2}
      \left[
      \arcsin\left(\frac{x}{a}\right)
      + \frac{x}{a}\sqrt{a^2-x^2}\,
      \right] + C
      & & \text{Justificativa em breve}}
  \end{align*}
\end{frame}

%------------------------------------------------------------------------------%
\begin{frame}{Transformação Inversa}
  Substituição simples -- a nova variável é função da original
  \[
    u = f(x)
  \]

  \pause
  Substituição inversa -- a variável original é função da nova
  \[
    x = g(u)
  \]
  \pause
  É necessário inverter $g$
\end{frame}

%------------------------------------------------------------------------------%
\begin{frame}{Substituições Trigonométricas}
\addtext{16}{0.35,2.2}{
\renewcommand*{\arraystretch}{2.3}
\begin{tabular}{c<{\;}>{\;}c<{\;}>{\;}c<{\;}>{\;}c} \hline
  Expressão
  & Substituição
  & Intervalo
  & Identidade \\ \hline
  \onslide<+->{\(\sqrt{a^2-x^2}\)
  & \(x=a\sin(\theta)\)
  & \(-\frac{\pi}{2} \leq \theta \leq \frac{\pi}{2}\)
  & \(1-\sin^2(\theta) = \cos^2(\theta)\)} \\
  \onslide<+->{\(\sqrt{a^2+x^2}\)
  & \(x=a\tan(\theta)\)
  & \(-\frac{\pi}{2}< \theta < \frac{\pi}{2} \)
  & \(1+\tan^2(\theta)= \sec^2(\theta)\)} \\
  \onslide<+->{\(\sqrt{x^2-a^2}\)
  & \(x=a\sec(\theta)\)
  & \(0 \leq \theta < \frac{\pi}{2} \) ou \(\pi \leq \theta < \frac{3\pi}{2} \)
  & \(\sec^2(\theta)-1= \tan^2(\theta)\)} \\ \hline
\end{tabular}}
\end{frame}

% 1
%------------------------------------------------------------------------------%
\addtocounter{exemplo}{1}
\section{Exemplo \theexemplo}

%------------------------------------------------------------------------------%
\begin{frame}{Exemplo \theexemplo}
  Calcule \(\int \frac{\sqrt{9-x^2}}{x^2}dx\)

  \pause
  \vspace{0.5\baselineskip}
  Note
  \(
    \sqrt{9-x^2} = \sqrt{3^2-x^2}
  \)

  \pause
  \vspace{\baselineskip}
  Substituição
  \[
    x = 3\sin(\theta)
    \qquad\qquad
    dx = 3 \cos(\theta) d\theta
    \qquad\text{com}\qquad
    -\frac{\pi}{2} \leq x \leq \frac{\pi}{2}
  \]
\end{frame}

%------------------------------------------------------------------------------%
\begin{frame}{Exemplo \theexemplo}
  \begin{align*}
    \onslide<+->{F & = \int \frac{\sqrt{9-x^2}}{x^2}dx          & & x = 3\sin(\theta)                              \\[2mm]}
    \onslide<+->{  & = \int \frac{\sqrt{9-9\sin^2(\theta)}}{9 \sin^2(\theta)} \, 3 \cos(\theta) d\theta            \\[2mm]}
    \onslide<+->{  & = \int \frac{3\times3}{9}\frac{\sqrt{1-\sin^2(\theta)}}{ \sin^2(\theta)} \cos(\theta) d\theta \\[2mm]}
    \onslide<+->{  & = \int \frac{\sqrt{1-\sin^2(\theta)}}{ \sin^2(\theta)} \cos(\theta) d\theta}
  \end{align*}
\end{frame}

%------------------------------------------------------------------------------%
\begin{frame}{Exemplo \theexemplo}
  \begin{align*}
    \onslide<+->{F & = \int \frac{\sqrt{1-\sin^2(\theta)}}{ \sin^2(\theta)} \cos(\theta) d\theta
                   & 1-\sin^2(\theta)= \cos^2(\theta) \\[2mm]}
    \onslide<+->{  & = \int \frac{\sqrt{\cos^2(\theta)}}{ \sin^2(\theta)} \cos(\theta) d\theta        \\[2mm]}
    \onslide<+->{  & = \int \frac{\cos^2(\theta)}{\sin^2(\theta)} d\theta                        \\[2mm]}
    \onslide<+->{  & = \int \cot^2(\theta) d\theta}
  \end{align*}
\end{frame}

%------------------------------------------------------------------------------%
\begin{frame}{Exemplo \theexemplo}
  \begin{align*}
    \onslide<+->{F & = \int \cot^2(\theta) d\theta       & \cot^2(\theta) = \csc^2(\theta)-1 \\[2mm]}
    \onslide<+->{  & = \int \csc^2(\theta) -1 \, d\theta & \big(\cot(\theta)\big)' = -\csc^2(\theta)\\[2mm]}
    \onslide<+->{  & = - \cot(\theta) - \theta +C        &}
  \end{align*}
  \onslide<+->{Precisamos voltar para a variável original \(x = 3\sin(\theta)\)}

  \onslide<+->{\[ \theta = \arcsin\left(\frac{x}{3}\right)\]}
  \onslide<+->{
  \[
    F = - {\color<7>{red}\cot\left(\arcsin\left(\frac{x}{3}\right)\right)}
        - \arcsin\left(\frac{x}{3}\right) + C
    \qquad
    \onslide<7>{\text{Forma ruim}}
  \]}
\end{frame}

%------------------------------------------------------------------------------%
\begin{frame}{Exemplo \theexemplo}
  \onslide<+->{\[
    x = 3\sin(\theta) \quad\Rightarrow\quad
    \sin(\theta)
    = \frac{x}{3}
    = \frac{\text{cateto oposto}}{\text{hipotenusa}}
  \]}

  \addgraphic{6}{10,3}{exemplo_triangulo_1}

  \onslide<+->{Por Pitágoras, o cateto adjacente é \(\sqrt{9-x^2}\)}

  \onslide<+->{Portanto
  \[
    \tan(\theta)
    = \frac{\text{cateto oposto}}{\text{cateto adjacente}}
  \]}

  \onslide<+->{\[
    \cot(\theta)
    = \frac{\text{cateto adjacente}}{\text{cateto oposto}}
    = \frac{\sqrt{9-x^2}}{x}
  \]}
\end{frame}

%------------------------------------------------------------------------------%
\begin{frame}{Exemplo \theexemplo}
  \begin{align*}
    \onslide<+->{\int \frac{\sqrt{9-x^2}}{x^2}dx
    & = -\cot(\theta) - \theta + C  \\[5mm]}
    \onslide<+->{& = -\frac{\sqrt{9-x^2}}{x} - \arcsin\left(\frac{x}{3}\right) + C}
  \end{align*}
\end{frame}

% 2
%------------------------------------------------------------------------------%
\addtocounter{exemplo}{1}
\section{Exemplo \theexemplo}

%------------------------------------------------------------------------------%
\begin{frame}{Exemplo \theexemplo}
  Encontre \(\int \frac{dx}{x^2\sqrt{x^2+4}}\)

  \pause
  \vspace{0.5\baselineskip}
  Note
  \(
    \sqrt{x^2+4} = \sqrt{2^2+x^2}
  \)

  \pause
  \vspace{\baselineskip}
  Substituição
  \[
    x = 2\tan(\theta)
    \qquad\qquad
    dx = 2\sec^2(\theta) d\theta
    \qquad\text{com}\qquad
    -\frac{\pi}{2} < \theta < \frac{\pi}{2}
  \]
\end{frame}

%------------------------------------------------------------------------------%
\begin{frame}{Exemplo \theexemplo}
  \begin{align*}
    \onslide<+->{F & = \int \frac{1}{x^2\sqrt{x^2+4}}dx
                   & & x = 2\tan(\theta)                       \\[2mm]}
    \onslide<+->{  & = \int \frac{1}{4\tan^2(\theta) \sqrt{4\tan^2(\theta)+4}}\,2\sec^2(\theta) d\theta      \\[2mm]}
    \onslide<+->{  & = \int \frac{\sec^2(\theta)}{2\tan^2(\theta) \sqrt{4\big(1+ \tan^2(\theta)\big)}} d\theta \\[2mm]}
    \onslide<+->{  & = \int \frac{\sec^2(\theta)}{4\tan^2(\theta) \sqrt{1+ \tan^2(\theta)}} d\theta
                   & & 1 + \tan^2(\theta) = \sec^2(\theta) \\[2mm]}
    \onslide<+->{  & = \frac{1}{4} \int \frac{\sec^2(\theta)}{\tan^2(\theta) \sqrt{\sec^2(\theta)}} d\theta}
    %
  \end{align*}
\end{frame}

%------------------------------------------------------------------------------%
\begin{frame}{Exemplo \theexemplo}
  \begin{align*}
    \onslide<+->{F & = \frac{1}{4} \int \frac{\sec^2(\theta)}{\tan^2(\theta) \sqrt{\sec^2(\theta)}} d\theta\\[3mm]}
%    & = \int \frac{\sec^2(\theta)}{\tan^2(\theta) \sec(\theta)}  d\theta \\[2mm]
    \onslide<+->{  & = \frac{1}{4} \int \frac{\sec(\theta)}{\tan^2(\theta)} d\theta              \\[3mm]}
    \onslide<+->{  & = \frac{1}{4} \int \frac{1}{\cos(\theta)}
                                   \left(\frac{\cos(\theta)}{\sin(\theta)}\right)^{\!\!2} d\theta \\[3mm]}
    \onslide<+->{  & = \frac{1}{4} \int \frac{\cos(\theta)}{\sin^2(\theta)} d\theta}
  \end{align*}
\end{frame}

%------------------------------------------------------------------------------%
\begin{frame}{Exemplo \theexemplo}
  Integrar
  \(
    F = \frac{1}{4} \int \frac{\cos(\theta)}{\sin^2(\theta)} d\theta
  \)

  \pause
  \vspace{0.5\baselineskip}
  Substituição
  \(
    u = \sin(\theta)
    \qquad
    du = \cos(\theta) d\theta
  \)

  \pause
  \vspace{0.5\baselineskip}
  Temos
  \[
    F
    = \frac{1}{4} \int \frac{du}{u^2}
    \pause
    = \frac{1}{4} \int u^{-2} du
    \pause
    = \frac{1}{4} \left(\frac{u^{-1}}{-1} + C\right)
    \pause
    = \frac{-1}{4\sin(\theta)} + C_1
  \]
\end{frame}

%------------------------------------------------------------------------------%
\begin{frame}{Exemplo \theexemplo}
  \[
    x= 2 \tan(\theta) \qquad\Rightarrow\qquad \tan(\theta) = \frac{x}{2}
  \]

  \addgraphic{6}{10,3}{exemplo_triangulo_2}

  \pause
  \[
    \frac{1}{\sin(\theta)}
    = \csc(\theta)
    \pause
    = \frac{\text{hipotenusa}}{\text{cateto oposto}}
    \pause
    = \frac{\sqrt{x^2+4}}{x}
  \]

  \pause
  \begin{align*}
    \onslide<+->{F & = \int \frac{dx}{x^2\sqrt{x^2+4}}  \\[2mm]}
    \onslide<+->{  & = \frac{-1}{4\sin(\theta)} + C \\[2mm]}
    \onslide<+->{  & = -\frac{\sqrt{x^2+4}}{4x} + C}
  \end{align*}
\end{frame}

% 3
%------------------------------------------------------------------------------%
\addtocounter{exemplo}{1}
\section{Exemplo \theexemplo}

%------------------------------------------------------------------------------%
\begin{frame}{Exemplo \theexemplo}
  Encontre \(\int\!\!\frac{dx}{\sqrt{x^2-a^2}}\) com $a>0$

  \pause
  \vspace{0.5\baselineskip}
  Substituição
  \[
    x = a\sec(\theta)
    \qquad
    dx = a \sec(\theta)\tan(\theta) d\theta
  \]
  onde \(0 < \theta < \frac{\pi}{2}\) ou \(\pi < \theta  < \frac{3\pi}{2}\)

  \pause
  \vspace{0.5\baselineskip}
  \[
    F = \int \frac{ a \sec(\theta) \tan(\theta) }
                  { \sqrt{a^2 \sec^2(\theta) - a^2 }} d\theta
  \]
\end{frame}

%------------------------------------------------------------------------------%
\begin{frame}{Exemplo \theexemplo}
\begin{columns}
\column{0.5\linewidth}
  \begin{align*}
    \onslide<+->{F & = \int \frac{dx}{\sqrt{x^2-a^2}}                                       \\[4mm]}
    \onslide<+->{  & = \int  \frac{ a \sec(\theta) \tan(\theta)}{\sqrt{a^2 \sec^2(\theta)-a^2}}  d\theta \\[4mm]}
    \onslide<+->{  & = \int \frac{a \sec(\theta) \tan(\theta)}{a\sqrt{\sec^2(\theta)-1}}  d\theta}
  \end{align*}
\column{0.5\linewidth}
  \onslide<+->{Usando \(\sec^2(\theta)-1= \tan^2(\theta)\)}
  \begin{align*}
    \onslide<+->{F & = \int \frac{\sec(\theta) \tan(\theta)}{\sqrt{\tan^2(\theta)}} d\theta \\[4mm]}
    \onslide<+->{  & = \int \frac{\sec(\theta)\tan(\theta)}{\tan(\theta)} d\theta           \\[4mm]}
    \onslide<+->{  & = \int \sec(\theta) d\theta                                            }
  \end{align*}
\end{columns}
\end{frame}

%------------------------------------------------------------------------------%
\begin{frame}{Exemplo \theexemplo}
  \begin{align*}
    \onslide<+->{F & = \int \sec(\theta) d\theta \\[2mm]}
    \onslide<+->{  & = \int \sec(\theta)
                          \;\frac{\sec(\theta)+\tan(\theta)}
                                 {\sec(\theta)+\tan(\theta)}\; d\theta \\[2mm]}
    \onslide<+->{  & = \int \frac{\sec^2(\theta)+ \sec(\theta)\tan(\theta)}
                                 {\sec(\theta)+\tan(\theta)} d\theta}
  \end{align*}

  \onslide<+->{Substituição
  \(
    u = \sec(\theta) + \tan(\theta)
    \qquad
    du = \big(\sec(\theta)\tan(\theta) + \sec^2(\theta)\big) d\theta
  \)}
  \[
    \onslide<+->{F = \int \frac{du}{u}}
    \onslide<+->{  = \ln\abs{u}+C}
    \onslide<+->{  = \ln \abs{\sec(\theta)+\tan(\theta)}+ C}
  \]
\end{frame}

%------------------------------------------------------------------------------%
\begin{frame}{Exemplo \theexemplo}
  \[
    \onslide<+->{x = a\sec(\theta)}
    \onslide<+->{\quad\Rightarrow\quad
    \sec(\theta)
      = \frac{\text{hipotenusa}}{\text{cateto adjacente}}
      = \frac{x}{a}}
  \]

  \addgraphic{6}{10,3}{exemplo_triangulo_3}

  \onslide<+->{\[
    \tan(\theta)
    = \frac{\text{cateto oposto}}{\text{cateto adjacente}}
    = \frac{\sqrt{x^2-a^2}}{a}
  \]}

  \begin{align*}
    \onslide<+->{F & = \int \frac{dx}{\sqrt{x^2-a^2}}                     \\[2mm]}
    \onslide<+->{  & = \ln \abs{\sec(\theta)+ \tan(\theta)} + C           \\[2mm]}
    \onslide<+->{  & = \ln \abs{\frac{x}{a}+ \frac{\sqrt{x^2-a^2}}{a}}+ C}
  \end{align*}
\end{frame}

% 4
%------------------------------------------------------------------------------%
\addtocounter{exemplo}{1}
\section{Exemplo \theexemplo}

%------------------------------------------------------------------------------%
\begin{frame}{Exemplo \theexemplo}
  \onslide<1->{Encontre \(\int \frac{x}{\sqrt{x^2+4}} dx\)}

  \vspace{0.5\baselineskip}
  \only<1|handout:0>{\vphantom{Substituição trigonométrica \(x = 2 \tan(\theta)\)}}
  \only<2|handout:0>{Substituição trigonométrica \(x = 2 \tan(\theta)\)}
  \only<3:->{\sout{Substituição trigonométrica \(x = 2 \tan(\theta)\)}}

  \vspace{0.5\baselineskip}
  \onslide<4->{Substituição simples \(u=x^2+4 \qquad du=2x dx\)}
  \[
    \onslide<5->{\int \frac{x}{\sqrt{x^2+4}} dx
                 = \frac{1}{2} \int \frac{du}{\sqrt{u}}}
    \onslide<6->{= \sqrt{u}+ C}
    \onslide<7->{= \sqrt{x^2+4}+C}
  \]
\end{frame}

% 5
%------------------------------------------------------------------------------%
\addtocounter{exemplo}{1}
\section{Exemplo \theexemplo}

%------------------------------------------------------------------------------%
\begin{frame}{Exemplo \theexemplo}
  Calcule \(\int x \arcsin(x) dx\)

  \pause
  \vspace{\baselineskip}
  Vários métodos serão necessários
\end{frame}

%------------------------------------------------------------------------------%
\begin{frame}{Exemplo \theexemplo\ -- Integral por Partes}
  \onslide<+->{\[
    F = \int x \arcsin(x) dx
  \]}
  \begin{align*}
    \onslide<+->{u  & = \arcsin(x)              & dv & = xdx           \\}
    \onslide<+->{du & =\frac{1}{\sqrt{1-x^2}}dx & v  & = \frac{x^2}{2}}
  \end{align*}

  \[
    \onslide<+->{F = uv - \int v du}
    \onslide<+->{  = \frac{x^2}{2} \arcsin(x)}
    \onslide<+->{  - \frac{1}{2}\int \frac{x^2}{\sqrt{1-x^2}}dx}
  \]
\end{frame}

%------------------------------------------------------------------------------%
\begin{frame}{Exemplo \theexemplo\ -- Substituição Trigonométrica}
  \onslide<+->{Resolver \( \int \frac{x^2}{\sqrt{1-x^2}}dx \)}

  \vspace{\baselineskip}
  \onslide<+->{Substituição trigonométrica
  \[
    x  = \sin(\theta)
    \qquad
    dx = \cos(\theta)d\theta
    \quad\text{com}\quad
    -\frac{\pi}{2} \leq \theta \leq \frac{\pi}{2}
  \]}
  \[
    \onslide<+->{G = \int \frac{x^2}{\sqrt{1-x^2}}dx}
    \onslide<+->{  = \int \frac{\sin^2(\theta)}{\sqrt{1-\sin^2(\theta)}} \cos(\theta)d\theta}
  \]
\end{frame}

%------------------------------------------------------------------------------%
\begin{frame}{Exemplo \theexemplo\ -- Substituição Trigonométrica}
  \begin{align*}
    \onslide<+->{G & = \int \frac{x^2}{\sqrt{1-x^2}}dx                                         \\[2mm]}
    \onslide<+->{  & = \int \frac{\sin^2(\theta)}{\sqrt{1-\sin^2(\theta)}} \cos(\theta)d\theta \\[2mm]}
    \onslide<+->{  & = \int \frac{\sin^2(\theta)}{\sqrt{\cos^2(\theta)}} \cos(\theta)d\theta   \\[2mm]}
    \onslide<+->{  & = \int \sin^2(\theta) d\theta}
  \end{align*}
\end{frame}

%------------------------------------------------------------------------------%
\begin{frame}{Exemplo \theexemplo\ -- Integral Trigonométrica}
  \onslide<+->{Calcular \( G = \int \sin^2(\theta) d\theta \)}

  \vspace{\baselineskip}
  \onslide<+->{Manipulação algébrica
  \[
    \sin^2(\theta)= \frac{1}{2}\big(1-\cos(2 \theta)\big)
  \]}
\end{frame}

%------------------------------------------------------------------------------%
\begin{frame}{Exemplo \theexemplo\ -- Integral Trigonométrica}
  \begin{align*}
    \onslide<+->{G & = \int \sin^2(\theta) d \theta                                            \\[2mm]}
    \onslide<+->{  & = \frac{1}{2}\int\big(1-\cos(2 \theta)\big) d\theta                       \\[2mm]}
    \onslide<+->{  & = \frac{1}{2}\left(\theta -\frac{\sin(2\theta)}{2}\right) + C             \\[2mm]}
    \onslide<+->{  & = \frac{1}{2}\left(\theta -\frac{2\sin(\theta)\cos(\theta)}{2}\right) + C \\[2mm]}
    \onslide<+->{  & = \frac{1}{2}\big(\theta -\sin(\theta)\cos(\theta)\big) + C}
  \end{align*}
\end{frame}

%------------------------------------------------------------------------------%
\begin{frame}[t]{Exemplo \theexemplo\ -- Desfazendo a Transformação Trigonométrica}
  \[
    \onslide<+->{ x
      = \sin(\theta)
      = \frac{\text{cateto oposto}}{\text{hipotenusa}}}
    \qquad
    \onslide<+->{
      \theta=\arcsin(x)}
  \]

  \addgraphic{4}{10,3}{exemplo_triangulo_4}

  \onslide<+->{\[
    \cos(\theta)= \sqrt{1-x^2}
  \]}
  \begin{align*}
    \onslide<+->{G & = \int \frac{x^2}{\sqrt{1-x^2}}dx                       \\[2mm]}
    \onslide<+->{  & = \frac{1}{2}\big(\theta -\sin(\theta)\cos(\theta)\big) \\[2mm]}
    \onslide<+->{  & = \frac{1}{2}\left( \arcsin(x)-x\sqrt{1-x^2} \right) + C}
  \end{align*}
\end{frame}

%------------------------------------------------------------------------------%
\begin{frame}{Exemplo \theexemplo\ -- Retornando a Integral Original}
  \begin{align*}
    \onslide<+->{F & = \int x \arcsin(x) dx                                 \\[5mm]}
    \onslide<+->{  & = \frac{x^2}{2} \arcsin(x)
                     - \frac{1}{2}\int \frac{x^2}{\sqrt{1-x^2}}dx           \\[5mm]}
    \onslide<+->{  & = \frac{x^2}{2} \arcsin(x)
                     - \frac{1}{4}\left(\arcsin(x)-x\sqrt{1-x^2}\right) + C \\[5mm]}
    \onslide<+->{  & = \left(\frac{x^2}{2} - \frac{1}{4}\right) \arcsin(x)
                     + \frac{x}{4} \sqrt{1-x^2}
                     + C}
  \end{align*}
\end{frame}

% 6
%------------------------------------------------------------------------------%
\addtocounter{exemplo}{1}
\section{Exemplo \theexemplo}

%------------------------------------------------------------------------------%
\begin{frame}{Exemplo \theexemplo}
  Calcule \(\int \frac{dt}{\sqrt{t^2-6t+13}}\)

  \pause
  \vspace{\baselineskip}
  Completando quadrado temos
  \[
    t^2 - 6t + 13
    \pause
    \;=\; (t-3)^2 - 9 + 13
    \pause
    \;=\; (t-3)^2 + 4
  \]

  \pause
  \[
    F = \int \frac{dt}{\sqrt{t^2-6t+13}}
    \pause
      = \int \frac{dt}{\sqrt{(t-3)^2 + 4}}
  \]
\end{frame}

%------------------------------------------------------------------------------%
\begin{frame}{Exemplo \theexemplo\ -- Substituição Simples}
  \onslide<+->{\[
    u = t - 3 \qquad\qquad du = dt
  \]}
  \vspace{-\baselineskip}
  \begin{align*}
    \onslide<+->{F & = \int \frac{dt}{\sqrt{(t-3)^2+4}}   \\[2mm]}
    \onslide<+->{  & = \int \frac{du}{\sqrt{u^2+4}}       \\[2mm]}
    \onslide<+->{  & = \int \frac{du}{\sqrt{u^2+2^2}}}
  \end{align*}
\end{frame}

%------------------------------------------------------------------------------%
\begin{frame}{Exemplo \theexemplo\ -- Substituição Trigonométrica}
  \onslide<+->{\[
    u  = 2\tan(\theta) \quad\Rightarrow\quad
    du = 2\sec^2(\theta) d\theta
  \]}
  \begin{align*}
    \onslide<+->{F & = \int \frac{du}{\sqrt{u^2+2^2}}                                  \\[2mm]}
    \onslide<+->{  & = \int \frac{2\sec^2(\theta) }{\sqrt{4\tan^2(\theta)+4}}d\theta}
  \end{align*}
\end{frame}

%------------------------------------------------------------------------------%
\begin{frame}{Exemplo \theexemplo\ -- Substituição Trigonométrica}
  \begin{align*}
    \onslide<+->{F & = \int \frac{2\sec^2(\theta) }{\sqrt{4\tan^2(\theta)+4}}d\theta   \\[2mm]}
    \onslide<+->{  & = \int \frac{2\sec^2(\theta) }{2\sqrt{\tan^2(\theta)+1}}d\theta
      & \tan^2(\theta)+1= \sec^2(\theta)                                  \\[2mm]}
    \onslide<+->{  & = \int \frac{\sec^2(\theta)}{\sec(\theta)}d\theta                 \\[2mm]}
    \onslide<+->{  & = \int \sec(\theta)d\theta                                        \\[2mm]}
    \onslide<+->{  & = \ln \abs{\sec(\theta)+ \tan(\theta)} + K}
  \end{align*}
\end{frame}

%------------------------------------------------------------------------------%
\begin{frame}{Exemplo \theexemplo\ -- Desfazendo as Transformações}
  \onslide<+->{\[
    u = 2\tan(\theta)
    \quad\Rightarrow\quad
    \tan(\theta) = \frac{u}{2} = \frac{t-3}{2}
  \]}

  \addgraphic{5.5}{10,3}{exemplo_triangulo_5}

  \onslide<+->{Hipotenusa}
  \begin{align*}
    \onslide<2->{\sqrt{(t-3)^2+2^2}
                 & = \sqrt{t^2-6t+9+4} \\}
    \onslide<+->{& = \sqrt{t^2-6t+13}}
  \end{align*}
  \onslide<+->{
  \[
    \sec(\theta)=\frac{\sqrt{t^2-6t+13}}{2}
  \]}
\end{frame}

%------------------------------------------------------------------------------%
\begin{frame}{Exemplo \theexemplo\ -- Desfazendo as Transformações}
  \begin{align*}
    \onslide<+->{F & = \int \frac{dt}{\sqrt{t^2-6t+13}}         \\[4mm]}
    \onslide<+->{  & = \ln \abs{\sec(\theta)+ \tan(\theta)} + K \\[4mm]}
    \onslide<+->{  & = \ln \abs{\frac{\sqrt{t^2-6t+13}}{2} + \frac{t-3}{2}} + K \\[4mm]}
    \onslide<+->{  & = \ln \abs{\sqrt{t^2-6t+13} + t-3} - \ln(2) + K \\[4mm]}
    \onslide<+->{  & = \ln \abs{\sqrt{t^2-6t+13} + t-3} + K_1}
  \end{align*}
\end{frame}

%------------------------------------------------------------------------------%
\begin{frame}
  \tableofcontents
\end{frame}

%------------------------------------------------------------------------------%
\end{document}
%------------------------------------------------------------------------------%